\documentclass[a4paper,twoside]{article}

\usepackage{morewrites}

\usepackage[svgnames,dvipsnames,x11names]{xcolor}
\usepackage{graphicx}
\usepackage[english]{babel}
\usepackage[colorlinks=true, linkcolor=NavyBlue, urlcolor=NavyBlue, citecolor=NavyBlue]{hyperref}
\usepackage[a4paper]{geometry}
\usepackage{ts-fonts}
\usepackage{amsmath}
\usepackage{float}
\usepackage{seqsplit}
\usepackage{ts-callouts}
\usepackage{ts-code}

\usepackage{lastpage}
\usepackage{refcount}
\makeatletter
\def\ps@plain{
  \let\@oddhead\@empty
  \let\@evenhead\@empty
  \def\@oddfoot{
    \hfil
    \ifnum\getpagerefnumber{LastPage}>1\relax
      \thepage
    \fi
    \hfil}
  \let\@evenfoot\@oddfoot
}
\makeatother

\title{Readers \& Writers}
\author{}\date{}

\begin{document}

\maketitle

\thispagestyle{plain}
TeXSmith converts a document through a typed intermediate representation (IR):

\begin{code}{text}{}{}
\begin{Verbatim}[commandchars=\\\{\},breaklines, breakanywhere, commandchars=\\\{\}]
read(HTML) → IR (texsmith.ir) → write(IR) → LaTeX
\end{Verbatim}
\end{code}
A \textbf{reader} lowers BeautifulSoup nodes into backend-agnostic IR nodes, and a
\textbf{writer} emits a backend (\LaTeX{}) from that IR. Custom constructs are added in
two halves that mirror how the built-in constructs work.

\section{Reader lowerings --- \texttt{@reads}}\label{reader-lowerings-texttt-reads}

A reader lowering is a callable decorated with \texttt{@reads(*tags, level, priority,
name)}. It \textbf{returns} an IR node (never mutates the tree). Returning
\texttt{NotHandled} lets the next candidate --- or the generic fallback --- run.

\begin{code}{python}{}{}
\begin{Verbatim}[commandchars=\\\{\},breaklines, breakanywhere, commandchars=\\\{\}]
\PY{k+kn}{from}\PY{+w}{ }\PY{n+nn}{bs4}\PY{+w}{ }\PY{k+kn}{import} \PY{n}{Tag}

\PY{k+kn}{from}\PY{+w}{ }\PY{n+nn}{texsmith}\PY{n+nn}{.}\PY{n+nn}{ir}\PY{+w}{ }\PY{k+kn}{import} \PY{n}{nodes} \PY{k}{as} \PY{n}{ir}
\PY{k+kn}{from}\PY{+w}{ }\PY{n+nn}{texsmith}\PY{n+nn}{.}\PY{n+nn}{readers}\PY{n+nn}{.}\PY{n+nn}{html}\PY{n+nn}{.}\PY{n+nn}{registry}\PY{+w}{ }\PY{k+kn}{import} \PY{n}{NotHandled}\PY{p}{,} \PY{n}{ReadLevel}\PY{p}{,} \PY{n}{reads}

\PY{n+nd}{@reads}\PY{p}{(}\PY{l+s+s2}{\PYZdq{}}\PY{l+s+s2}{span}\PY{l+s+s2}{\PYZdq{}}\PY{p}{,} \PY{n}{level}\PY{o}{=}\PY{n}{ReadLevel}\PY{o}{.}\PY{n}{INLINE}\PY{p}{,} \PY{n}{name}\PY{o}{=}\PY{l+s+s2}{\PYZdq{}}\PY{l+s+s2}{data\PYZus{}counter}\PY{l+s+s2}{\PYZdq{}}\PY{p}{,} \PY{n}{priority}\PY{o}{=}\PY{l+m+mi}{50}\PY{p}{)}
\PY{k}{def}\PY{+w}{ }\PY{n+nf}{read\PYZus{}data\PYZus{}counter}\PY{p}{(}\PY{n}{tag}\PY{p}{:} \PY{n}{Tag}\PY{p}{,} \PY{n}{ctx}\PY{p}{)} \PY{o}{\PYZhy{}}\PY{o}{\PYZgt{}} \PY{n}{ir}\PY{o}{.}\PY{n}{Span} \PY{o}{|} \PY{n+nb}{object}\PY{p}{:}
    \PY{n}{classes} \PY{o}{=} \PY{n}{tag}\PY{o}{.}\PY{n}{get}\PY{p}{(}\PY{l+s+s2}{\PYZdq{}}\PY{l+s+s2}{class}\PY{l+s+s2}{\PYZdq{}}\PY{p}{)} \PY{o+ow}{or} \PY{p}{[}\PY{p}{]}
    \PY{n}{tokens} \PY{o}{=} \PY{p}{\PYZob{}}\PY{n}{classes}\PY{p}{\PYZcb{}} \PY{k}{if} \PY{n+nb}{isinstance}\PY{p}{(}\PY{n}{classes}\PY{p}{,} \PY{n+nb}{str}\PY{p}{)} \PY{k}{else} \PY{n+nb}{set}\PY{p}{(}\PY{n}{classes}\PY{p}{)}
    \PY{k}{if} \PY{l+s+s2}{\PYZdq{}}\PY{l+s+s2}{data\PYZhy{}counter}\PY{l+s+s2}{\PYZdq{}} \PY{o+ow}{not} \PY{o+ow}{in} \PY{n}{tokens}\PY{p}{:}
        \PY{k}{return} \PY{n}{NotHandled}
    \PY{k}{return} \PY{n}{ir}\PY{o}{.}\PY{n}{Span}\PY{p}{(}\PY{n}{content}\PY{o}{=}\PY{p}{(}\PY{p}{)}\PY{p}{,} \PY{n}{attrs}\PY{o}{=}\PY{p}{(}\PY{p}{(}\PY{l+s+s2}{\PYZdq{}}\PY{l+s+s2}{role}\PY{l+s+s2}{\PYZdq{}}\PY{p}{,} \PY{l+s+s2}{\PYZdq{}}\PY{l+s+s2}{counter}\PY{l+s+s2}{\PYZdq{}}\PY{p}{)}\PY{p}{,}\PY{p}{)}\PY{p}{)}
\end{Verbatim}
\end{code}
\section{Writer emitters --- \texttt{@writes}}\label{writer-emitters-texttt-writes}

A writer emitter is a method decorated with \texttt{@writes(NodeType)} on a
\texttt{LaTeXWriter} subclass. Dispatch is typed by node class via the MRO; a node
without an emitter raises a clear \texttt{LaTeXWriteError}.

\begin{code}{python}{}{}
\begin{Verbatim}[commandchars=\\\{\},breaklines, breakanywhere, commandchars=\\\{\}]
\PY{k+kn}{from}\PY{+w}{ }\PY{n+nn}{texsmith}\PY{n+nn}{.}\PY{n+nn}{ir}\PY{+w}{ }\PY{k+kn}{import} \PY{n}{nodes} \PY{k}{as} \PY{n}{ir}
\PY{k+kn}{from}\PY{+w}{ }\PY{n+nn}{texsmith}\PY{n+nn}{.}\PY{n+nn}{writers}\PY{n+nn}{.}\PY{n+nn}{latex}\PY{+w}{ }\PY{k+kn}{import} \PY{n}{LaTeXWriter}\PY{p}{,} \PY{n}{writes}

\PY{k}{class}\PY{+w}{ }\PY{n+nc}{CountingWriter}\PY{p}{(}\PY{n}{LaTeXWriter}\PY{p}{)}\PY{p}{:}
    \PY{n+nd}{@writes}\PY{p}{(}\PY{n}{ir}\PY{o}{.}\PY{n}{Span}\PY{p}{)}
    \PY{k}{def}\PY{+w}{ }\PY{n+nf}{\PYZus{}counter\PYZus{}span}\PY{p}{(}\PY{n+nb+bp}{self}\PY{p}{,} \PY{n}{node}\PY{p}{:} \PY{n}{ir}\PY{o}{.}\PY{n}{Span}\PY{p}{)} \PY{o}{\PYZhy{}}\PY{o}{\PYZgt{}} \PY{n+nb}{str}\PY{p}{:}
        \PY{k}{if} \PY{n+nb}{dict}\PY{p}{(}\PY{n}{node}\PY{o}{.}\PY{n}{attrs}\PY{p}{)}\PY{o}{.}\PY{n}{get}\PY{p}{(}\PY{l+s+s2}{\PYZdq{}}\PY{l+s+s2}{role}\PY{l+s+s2}{\PYZdq{}}\PY{p}{)} \PY{o}{==} \PY{l+s+s2}{\PYZdq{}}\PY{l+s+s2}{counter}\PY{l+s+s2}{\PYZdq{}}\PY{p}{:}
            \PY{n}{value} \PY{o}{=} \PY{n+nb+bp}{self}\PY{o}{.}\PY{n}{state}\PY{o}{.}\PY{n}{state}\PY{o}{.}\PY{n}{next\PYZus{}counter}\PY{p}{(}\PY{l+s+s2}{\PYZdq{}}\PY{l+s+s2}{data\PYZhy{}counter}\PY{l+s+s2}{\PYZdq{}}\PY{p}{)}
            \PY{k}{return} \PY{l+s+sa}{f}\PY{l+s+s2}{\PYZdq{}}\PY{l+s+se}{\PYZbs{}\PYZbs{}}\PY{l+s+s2}{counter}\PY{l+s+se}{\PYZob{}\PYZob{}}\PY{l+s+si}{\PYZob{}}\PY{n}{value}\PY{l+s+si}{\PYZcb{}}\PY{l+s+se}{\PYZcb{}\PYZcb{}}\PY{l+s+s2}{\PYZdq{}}
        \PY{k}{return} \PY{n+nb}{super}\PY{p}{(}\PY{p}{)}\PY{o}{.}\PY{n}{\PYZus{}span}\PY{p}{(}\PY{n}{node}\PY{p}{)}
\end{Verbatim}
\end{code}
See \href{https://github.com/texsmith/texsmith/blob/main/examples/custom-render/counter.py}{\texttt{examples/custom-\allowbreak{}render/counter.py}}
for a complete, runnable extension wiring both halves into a \texttt{LaTeXRenderer}.

\begin{callout}[callout tip]{Tip}
Keep the IR semantic and backend-neutral: encode hints via \texttt{Span}/\texttt{Div}
\texttt{attrs} (e.g. \texttt{("role", "counter")}) rather than backend strings. Only the
writer knows about \LaTeX{}.
\end{callout}
\section{Shipping readers \& writers in a template}\label{shipping-readers-writers-in-a-template}

When custom constructs belong to a specific \textbf{template} (an exam, a thesis, a
poster...), declare the reader modules and the writer subclass in that template's
\texttt{[latex.template]} manifest section. TeXSmith resolves them when the template is
selected and applies them to the renderer \textbf{for that template only} --- other
templates keep the default bundled read\(\rightarrow\)write path.

\begin{code}{toml}{}{}
\begin{Verbatim}[commandchars=\\\{\},breaklines, breakanywhere, commandchars=\\\{\}]
\PY{k}{[}\PY{k}{latex}\PY{k}{.}\PY{k}{template}\PY{k}{]}
\PY{n}{name}\PY{+w}{ }\PY{o}{=}\PY{+w}{ }\PY{l+s+s2}{\PYZdq{}}\PY{l+s+s2}{exam}\PY{l+s+s2}{\PYZdq{}}
\PY{n}{version}\PY{+w}{ }\PY{o}{=}\PY{+w}{ }\PY{l+s+s2}{\PYZdq{}}\PY{l+s+s2}{1.0.0}\PY{l+s+s2}{\PYZdq{}}
\PY{n}{entrypoint}\PY{+w}{ }\PY{o}{=}\PY{+w}{ }\PY{l+s+s2}{\PYZdq{}}\PY{l+s+s2}{template/template.tex}\PY{l+s+s2}{\PYZdq{}}
\PY{n}{engine}\PY{+w}{ }\PY{o}{=}\PY{+w}{ }\PY{l+s+s2}{\PYZdq{}}\PY{l+s+s2}{lualatex}\PY{l+s+s2}{\PYZdq{}}

\PY{c+c1}{\PYZsh{} Modules whose @reads lowerings are layered on top of the bundled registry.}
\PY{c+c1}{\PYZsh{} A higher\PYZhy{}priority handler may return NotHandled to fall through to a core one.}
\PY{n}{readers}\PY{+w}{ }\PY{o}{=}\PY{+w}{ }\PY{p}{[}\PY{l+s+s2}{\PYZdq{}}\PY{l+s+s2}{my\PYZus{}exam\PYZus{}pkg.reader}\PY{l+s+s2}{\PYZdq{}}\PY{p}{]}

\PY{c+c1}{\PYZsh{} A \PYZdq{}module:Class\PYZdq{} reference to a LaTeXWriter subclass adding/overriding @writes.}
\PY{n}{writer}\PY{+w}{ }\PY{o}{=}\PY{+w}{ }\PY{l+s+s2}{\PYZdq{}}\PY{l+s+s2}{my\PYZus{}exam\PYZus{}pkg.writer:ExamLaTeXWriter}\PY{l+s+s2}{\PYZdq{}}
\end{Verbatim}
\end{code}
\begin{itemize}
\item{} \textbf{\texttt{readers}} --- a list of importable module paths. Every \texttt{@reads}-decorated
  callable found in each module is registered, layered on top of the bundled
  HTML\(\rightarrow\)IR lowerings (\texttt{texsmith.readers.html.build\_reader\_registry}).
\item{} \textbf{\texttt{writer}} --- a \texttt{"module:Class"} reference to a \texttt{LaTeXWriter} subclass. It is
  validated at render time (must subclass \texttt{LaTeXWriter}), then used as the
  template's writer so its \texttt{@writes} emitters and overrides take effect.

\end{itemize}
Both are resolved with the same import machinery as attribute normalisers and
fragment entrypoints, so a typo or a bad reference fails early with an
actionable \texttt{TemplateError}. A template that declares neither is unaffected.

\begin{callout}[callout note]{Scope}
These hooks are \textbf{template-scoped by design}. Because exam-style rules
(e.g. mapping every \texttt{h1} to \texttt{\textbackslash{}question}) would corrupt unrelated documents,
they are never registered globally --- only while the declaring template
renders. There is no global reader/writer entry point.
\end{callout}
Inside an extension, read template front-matter overrides from
\texttt{self.state.runtime["template\_overrides"]} (already populated by the pipeline);
keep cross-node state on \texttt{self.state.state} (the \texttt{DocumentState}) and harvest it
in a pre-pass over the IR (\texttt{texsmith.ir.visitor.walk}) rather than mutating
shared state during emission.

\section{Reference}\label{reference}

::: texsmith.readers.html.registry

::: texsmith.readers.html.reader

::: texsmith.writers.latex

\end{document}
